When six metrics are presented simultaneously, courts may ask the expert witness
to rank them or explain apparent conflicts.
This section provides guidance on weighting and conflict resolution.

\subsection{When Metrics Agree}

When all six metrics classify a district as compact (PP $\geq 0.20$,
Reock $\geq 0.35$, CH $\geq 0.88$, S $\leq 2.2$, LW $\leq 2.5$, PWC within
two standard deviations), the composite profile provides an unambiguous conclusion:
the district is compact by all standard criteria.
This is the strongest possible evidentiary posture for an expert witness.

\subsection{When Metrics Disagree}

Metric disagreement arises in predictable patterns:

\paragraph{High PP, low Reock.}
Occurs for districts with smooth perimeters (high PP) but elongated shapes
whose MBC is large (low Reock).
Interpretation: the district is not jagged, but it is elongated.
Lead with LW to characterise the elongation; note that PP's higher value
reflects good boundary quality, not good overall shape.

\paragraph{High Reock, low CH.}
Occurs for districts with tentacle arms (low CH) whose extreme points determine
a moderate MBC (moderate Reock).
Interpretation: the district has appendages. Lead with CH as the diagnostic metric.

\paragraph{Low PWC, high geometric metrics.}
Occurs for compact urban districts where dense population is concentrated
near the centre (low PWC) but the boundary is irregular (lower PP, Reock).
Interpretation: the district serves its population well despite irregular shape.
PWC's low value provides a positive counter-narrative.

\subsection{Recommended Weighting}

For litigation purposes, we recommend equal weighting of the five geometric
metrics (PP, Reock, CH, S, LW) and presenting PWC separately as the
representational dimension.
This avoids the criticism that the expert is applying arbitrary weights to
favour a particular metric.
The composite verdict is:
\begin{enumerate}
  \item If $\geq 4$ of 5 geometric metrics are in the ``acceptable'' range
    (K.7, Section~7), the district is compact.
  \item If $\leq 2$ of 5 geometric metrics are in the acceptable range,
    the district is non-compact.
  \item For intermediate cases ($3$ of $5$), the expert should characterise
    the specific dimension of non-compactness (elongation, tentacle, etc.)
    and explain its geographic cause.
\end{enumerate}
